\clearpage
\begin{explanationbox}
\textbf{\textcolor{CExplanation}{为什么要有三种写法？先看一道最简单的题}}
\begin{center}\begin{tikzpicture}[font=\small,>=stealth,line join=round]
\draw[gray!40,rounded corners=3pt] (0,0) rectangle (14.2,8.3);\node at (7.1,7.700000000000001) {三种写法，对应同一条抛物线};
\begin{scope}[shift={(1.7000000000000002,2.9642857142857144)},x=1.2cm,y=0.6571428571428571cm]\draw[->,gray] (-0.5,0)--(4.5,0) node[right]{$x$};\draw[->,gray] (0,-2)--(0,5) node[above]{$y$};\draw[blue,thick,solid,domain=-0.2:4.2,samples=180] plot(\x,{(\x-2)^2-1});\fill[orange!85!black] (0,3) circle(2.2pt);\fill[orange!85!black] (1,0) circle(2.2pt);\fill[orange!85!black] (3,0) circle(2.2pt);\fill[orange!85!black] (2,-1) circle(2.2pt);\draw[violet,dashed] (2,-.95)--(2,4.5);\node[left] at (0,3){$(0,3)$};\node[below=5pt] at (1,0){$1$};\node[below=5pt] at (3,0){$3$};\node[below=4pt] at (2,-1){$(2,-1)$};\end{scope}\node[anchor=west,align=left,text width=5.0cm] at (8.4,4.0) {$x^2-4x+3$\\[9pt]$=(x-2)^2-1$\\[9pt]$=(x-1)(x-3)$\\[9pt]顶点 $(2,-1)$\\[9pt]交点 $(1,0),(3,0)$};
\node[text width=12.6cm,align=center] at (7.1,.65) {一般式读纵轴交点，顶点式读顶点，两根式读横轴交点。};
\end{tikzpicture}\end{center}


假设你要画出一条抛物线，题目可能告诉你：
“它经过 $(0,3)$”；“最低点在 $(2,-1)$”；或者“它经过 $(1,0)$ 和 $(3,0)$”。
这些条件指向的是同一条抛物线的\textbf{不同位置}。
就像介绍一间教室，可以说门牌号、教室中心在哪里，或教室两扇门分别在哪里；不是三个教室，只是三种描述。

\medskip

\end{explanationbox}
\clearpage
\begin{explanationbox}
\textbf{\textcolor{CExplanation}{从一个具体例子，把三种写法连起来}}
\begin{center}\begin{tikzpicture}[font=\small,>=stealth,line join=round]
\draw[gray!40,rounded corners=3pt] (0,0) rectangle (14.2,8.3);\node at (7.1,7.700000000000001) {三种写法，对应同一条抛物线};
\begin{scope}[shift={(1.7000000000000002,2.9642857142857144)},x=1.2cm,y=0.6571428571428571cm]\draw[->,gray] (-0.5,0)--(4.5,0) node[right]{$x$};\draw[->,gray] (0,-2)--(0,5) node[above]{$y$};\draw[blue,thick,solid,domain=-0.2:4.2,samples=180] plot(\x,{(\x-2)^2-1});\fill[orange!85!black] (0,3) circle(2.2pt);\fill[orange!85!black] (1,0) circle(2.2pt);\fill[orange!85!black] (3,0) circle(2.2pt);\fill[orange!85!black] (2,-1) circle(2.2pt);\draw[violet,dashed] (2,-.95)--(2,4.5);\node[left] at (0,3){$(0,3)$};\node[below=5pt] at (1,0){$1$};\node[below=5pt] at (3,0){$3$};\node[below=4pt] at (2,-1){$(2,-1)$};\end{scope}\node[anchor=west,align=left,text width=5.0cm] at (8.4,4.0) {$x^2-4x+3$\\[9pt]$=(x-2)^2-1$\\[9pt]$=(x-1)(x-3)$\\[9pt]顶点 $(2,-1)$\\[9pt]交点 $(1,0),(3,0)$};
\node[text width=12.6cm,align=center] at (7.1,.65) {一般式读纵轴交点，顶点式读顶点，两根式读横轴交点。};
\end{tikzpicture}\end{center}


\[
\underbrace{x^2-4x+3}_{\text{一般式}}
=\underbrace{(x-2)^2-1}_{\text{顶点式}}
=\underbrace{(x-1)(x-3)}_{\text{两根式}}.
\]
这三个式子不只是“长得像”，无论 $x$ 取什么实数，算出的结果都一样。
\begin{itemize}[leftmargin=*,itemsep=0.35em,topsep=0.3em]
\item \textbf{看一般式：}当 $x=0$ 时，$y=3$，因此经过 $(0,3)$。
\item \textbf{看顶点式：}$(x-2)^2\ge0$，所以 $y=(x-2)^2-1\ge-1$；在 $x=2$ 时取到最小值 $-1$。
\item \textbf{看两根式：}只要 $x=1$ 或 $x=3$，就有一个括号等于 $0$，于是 $y=0$。
\end{itemize}


\noindent\small\textbf{读图：}图上只画了一条抛物线。橙点标出两个与 $x$ 轴的交点、最低点和与 $y$ 轴的交点，虚线是对称轴。\normalsize

\medskip

\end{explanationbox}
\clearpage
\begin{explanationbox}
\textbf{\textcolor{CExplanation}{配方为什么能找到顶点？}}
\begin{center}\begin{tikzpicture}[font=\small,>=stealth,line join=round]
\draw[gray!40,rounded corners=3pt] (0,0) rectangle (14.2,8.3);\node at (7.1,7.700000000000001) {配方：离轴越远，平方项越大};
\begin{scope}[shift={(2.6,2.9642857142857144)},x=1.5cm,y=0.6571428571428571cm]\draw[->,gray] (-1,0)--(3,0) node[right]{$x$};\draw[->,gray] (0,-2)--(0,5) node[above]{$y$};\draw[blue,thick,solid,domain=-0.5:2.5,samples=180] plot(\x,{2*(\x-1)^2-1});\fill[orange!85!black] (1,-1) circle(2.2pt);\fill[orange!85!black] (0,1) circle(2.2pt);\fill[orange!85!black] (2,1) circle(2.2pt);\draw[violet,dashed] (1,-1.8)--(1,4.6);\draw[orange!85!black,dashed] (0,1)--(2,1);\foreach \v in {0,1,2}{\node[below=4pt] at (\v,-1.7){$\v$};}\end{scope}\node[anchor=west,align=left,text width=5.0cm] at (8.4,4.0) {$2x^2-4x+1$\\[9pt]$=2(x-1)^2-1$\\[9pt]$x=1$：最小值 $-1$\\[9pt]$x=0,2$：高度都为1\\[9pt]到轴距离都为1};
\node[text width=12.6cm,align=center] at (7.1,.65) {平方项只看到轴的距离；配方把顶点位置和高度直接显露出来。};
\end{tikzpicture}\end{center}


先别急着记公式，从 $x^2-4x+3$ 出发：
\[
\begin{aligned}
x^2-4x+3
&=\bigl(x^2-4x+4\bigr)-4+3\\
&=(x-2)^2-1.
\end{aligned}
\]
为什么要“加 $4$ 再减 $4$”？因为 $x^2-4x+4$ 恰好是 $(x-2)^2$。
\textbf{把难看的一串式子变成“平方＋一个数”，就能直接判断它最小能有多小。}

\medskip

\end{explanationbox}
\clearpage
\begin{explanationbox}
\textbf{\textcolor{CExplanation}{两根式为什么要写两个括号？}}
\begin{center}\begin{tikzpicture}[font=\small,>=stealth,line join=round]
\draw[gray!40,rounded corners=3pt] (0,0) rectangle (14.2,8.3);\node at (7.1,7.700000000000001) {零点相同，a不同，曲线的高度与方向仍会不同};
\begin{scope}[shift={(1.7000000000000002,3.6214285714285714)},x=1.2cm,y=0.6571428571428571cm]\draw[->,gray] (-0.5,0)--(4.5,0) node[right]{$x$};\draw[->,gray] (0,-3)--(0,4) node[above]{$y$};\draw[blue,thick,solid,domain=0.4:3.6,samples=180] plot(\x,{(\x-1)*(\x-3)});\draw[teal,thick,solid,domain=0.4:3.6,samples=180] plot(\x,{2*(\x-1)*(\x-3)});\draw[violet,thick,solid,domain=0.4:3.6,samples=180] plot(\x,{-(\x-1)*(\x-3)});\fill[orange!85!black] (1,0) circle(2.2pt);\fill[orange!85!black] (3,0) circle(2.2pt);\fill[orange!85!black] (2,-1) circle(2.2pt);\fill[orange!85!black] (2,-2) circle(2.2pt);\fill[orange!85!black] (2,1) circle(2.2pt);\foreach \v in {1,2,3}{\node[below=4pt] at (\v,0){$\v$};}\end{scope}\node[anchor=west,align=left,text width=5.0cm] at (8.4,4.0) {蓝：$a=1$，顶点高 $-1$\\[9pt]青：$a=2$，顶点高 $-2$\\[9pt]紫：$a=-1$，顶点高 $1$\\[9pt]零点都为 $1,3$\\[9pt]还须额外条件确定a};
\node[text width=12.6cm,align=center] at (7.1,.65) {同尺度：两根式中的a控制开口，不能从两个零点中省去。};
\end{tikzpicture}\end{center}


小学、初中就学过：两个数相乘等于 $0$，至少有一个数等于 $0$。
因此要让函数在 $x=1$ 和 $x=3$ 时都等于 $0$，最自然的写法就是
\[
y=a(x-1)(x-3).
\]
前面的 $a$ 控制图象向上还是向下、开口显得宽还是窄。
只知道两个交点，\textbf{还不能确定 $a$}；还需要第三个条件，如经过某点、最高点或最低点的高度。

\medskip

\end{explanationbox}
\clearpage
\begin{explanationbox}
\textbf{\textcolor{CExplanation}{眼前一亮（一）：为什么两个交点的中点就是对称轴？}}
\begin{center}\begin{tikzpicture}[font=\small,>=stealth,line join=round]
\draw[gray!40,rounded corners=3pt] (0,0) rectangle (14.2,8.3);\node at (7.1,7.700000000000001) {平方差：两个零点关于中点等距};
\begin{scope}[shift={(1.5285714285714287,3.9499999999999997)},x=0.8571428571428571cm,y=0.7666666666666666cm]\draw[->,gray] (-0.5,0)--(6.5,0) node[right]{$x$};\draw[->,gray] (0,-3)--(0,3) node[above]{$y$};\draw[blue,thick,solid,domain=0.2:5.8,samples=180] plot(\x,{.5*(\x-1)*(\x-5)});\fill[orange!85!black] (1,0) circle(2.2pt);\fill[orange!85!black] (5,0) circle(2.2pt);\fill[orange!85!black] (3,-2) circle(2.2pt);\draw[violet,dashed] (3,-2.6)--(3,2.6);\draw[<->,orange!85!black] (1,1)--(3,1) node[midway,above]{$d=2$};\draw[<->,orange!85!black] (3,1)--(5,1) node[midway,above]{$d=2$};\foreach \v/\L in {1/{x_1=1},3/{m=3},5/{x_2=5}}{\node[below=4pt] at (\v,-2.7){$\L$};}\end{scope}\node[anchor=west,align=left,text width=5.0cm] at (8.4,4.0) {图取 $a=\frac12$\\[9pt]$x_1=1,\ x_2=5$\\[9pt]$m=3,\ d=2$\\[9pt]$f=\frac12[(x-3)^2-4]$\\[9pt]顶点 $(3,-2)$};
\node[text width=12.6cm,align=center] at (7.1,.65) {两根分别是m减d与m加d；将两括号相乘，就出现以m为轴的平方项。};
\end{tikzpicture}\end{center}


把两根式写成
\[
y=a(x-x_1)(x-x_2).
\]
令 $m=\frac{x_1+x_2}{2}$，$d=\frac{x_2-x_1}{2}$，于是 $x_1=m-d$、$x_2=m+d$。
所以
\[
\begin{aligned}
y&=a(x-m+d)(x-m-d)\\
 &=a\bigl[(x-m)^2-d^2\bigr].
\end{aligned}
\]
现在看清楚了：它就是以 $x=m$ 为对称轴的顶点式。两个交点向左、向右离中点一样远，顶点正好在中间。
例如交点横坐标为 $1,3$，对称轴\textbf{不必列方程}，就是 $x=2$。

\medskip

\end{explanationbox}
\clearpage
\begin{explanationbox}
\textbf{\textcolor{CExplanation}{眼前一亮（二）：判别式其实还告诉你顶点离 $x$ 轴的关系}}
\begin{center}\begin{tikzpicture}[font=\small,>=stealth,line join=round]
\draw[gray!40,rounded corners=3pt] (0,0) rectangle (14.2,8.5);\node at (7.1,7.9) {开口向上：顶点高低决定实数交点个数};
\node at (2.4,6.9) {$\Delta=4>0$};\begin{scope}[shift={(2.4,3.8)},x=0.85cm,y=0.6cm]\draw[->,gray] (-2,0)--(2,0) node[right]{$x$};\draw[->,gray] (0,-2)--(0,4) node[above]{$y$};\draw[blue,thick,solid,domain=-1.7:1.7,samples=160] plot(\x,{(\x)^2-1});\fill[orange!85!black] (-1,0) circle(2pt);\fill[orange!85!black] (1,0) circle(2pt);\fill[orange!85!black] (0,-1) circle(2pt);\end{scope}\node[align=center,text width=4cm] at (2.4,1.7) {$y=x^2-1$\\两个交点};\node at (7.050000000000001,6.9) {$\Delta=0$};\begin{scope}[shift={(7.050000000000001,3.8)},x=0.85cm,y=0.6cm]\draw[->,gray] (-2,0)--(2,0) node[right]{$x$};\draw[->,gray] (0,-2)--(0,4) node[above]{$y$};\draw[teal,thick,solid,domain=-1.7:1.7,samples=160] plot(\x,{(\x)^2});\fill[orange!85!black] (0,0) circle(2pt);\end{scope}\node[align=center,text width=4cm] at (7.050000000000001,1.7) {$y=x^2$\\一个交点};\node at (11.7,6.9) {$\Delta=-4<0$};\begin{scope}[shift={(11.7,3.8)},x=0.85cm,y=0.6cm]\draw[->,gray] (-2,0)--(2,0) node[right]{$x$};\draw[->,gray] (0,-2)--(0,4) node[above]{$y$};\draw[violet,thick,solid,domain=-1.7:1.7,samples=160] plot(\x,{(\x)^2+1});\fill[orange!85!black] (0,1) circle(2pt);\end{scope}\node[align=center,text width=4cm] at (11.7,1.7) {$y=x^2+1$\\没有交点};
\node[text width=12.6cm,align=center] at (7.1,.65) {同尺度，$a=1$：$k=-1,0,1$，恒有$\Delta=-4ak$。};
\end{tikzpicture}\end{center}


因为
\[
k=c-\frac{b^2}{4a}=-\frac{b^2-4ac}{4a},
\qquad\boxed{\Delta=-4ak},
\]
所以，判断是否与 $x$ 轴相交，可以不必每次都重新解方程：
\begin{itemize}[leftmargin=*,itemsep=0.3em,topsep=0.25em]
\item 开口向上（$a>0$）：最低点在 $x$ 轴下方，交两次；正好在轴上，碰一次；在轴上方，不相交。
\item 开口向下（$a<0$）：最高点在 $x$ 轴上方，交两次；正好在轴上，碰一次；在轴下方，不相交。
\end{itemize}

\end{explanationbox}\clearpage\begin{explanationbox}
\textbf{\textcolor{CExplanation}{开口向下时，也要一起看a与k}}
\begin{center}\begin{tikzpicture}[font=\small,>=stealth,line join=round]
\draw[gray!40,rounded corners=3pt] (0,0) rectangle (14.2,8.5);\node at (7.1,7.9) {开口向下：最高点也要连同a的符号判断};
\node at (2.4,6.9) {$\Delta=4>0$};\begin{scope}[shift={(2.4,5.0)},x=0.85cm,y=0.6cm]\draw[->,gray] (-2,0)--(2,0) node[right]{$x$};\draw[->,gray] (0,-4)--(0,2) node[above]{$y$};\draw[blue,thick,solid,domain=-1.7:1.7,samples=160] plot(\x,{-(\x)^2+1});\fill[orange!85!black] (-1,0) circle(2pt);\fill[orange!85!black] (1,0) circle(2pt);\fill[orange!85!black] (0,1) circle(2pt);\end{scope}\node[align=center,text width=4cm] at (2.4,1.7) {$y=-x^2+1$\\两个交点};\node at (7.050000000000001,6.9) {$\Delta=0$};\begin{scope}[shift={(7.050000000000001,5.0)},x=0.85cm,y=0.6cm]\draw[->,gray] (-2,0)--(2,0) node[right]{$x$};\draw[->,gray] (0,-4)--(0,2) node[above]{$y$};\draw[teal,thick,solid,domain=-1.7:1.7,samples=160] plot(\x,{-(\x)^2});\fill[orange!85!black] (0,0) circle(2pt);\end{scope}\node[align=center,text width=4cm] at (7.050000000000001,1.7) {$y=-x^2$\\一个交点};\node at (11.7,6.9) {$\Delta=-4<0$};\begin{scope}[shift={(11.7,5.0)},x=0.85cm,y=0.6cm]\draw[->,gray] (-2,0)--(2,0) node[right]{$x$};\draw[->,gray] (0,-4)--(0,2) node[above]{$y$};\draw[violet,thick,solid,domain=-1.7:1.7,samples=160] plot(\x,{-(\x)^2-1});\fill[orange!85!black] (0,-1) circle(2pt);\end{scope}\node[align=center,text width=4cm] at (11.7,1.7) {$y=-x^2-1$\\没有交点};
\node[text width=12.6cm,align=center] at (7.1,.65) {同尺度，$a=-1$：$k=1,0,-1$，仍有$\Delta=-4ak$。};
\end{tikzpicture}\end{center}

\textbf{别单独看 $k$ 的正负；要连同开口方向 $a$ 一起看。}

\medskip

\end{explanationbox}
\clearpage
\begin{explanationbox}
\textbf{\textcolor{CExplanation}{眼前一亮（三）：从顶点式直接“读出”交点}}
\begin{center}\begin{tikzpicture}[font=\small,>=stealth,line join=round]
\draw[gray!40,rounded corners=3pt] (0,0) rectangle (14.2,8.3);\node at (7.1,7.700000000000001) {例1：顶点、对称轴和两个零点同时读出};
\begin{scope}[shift={(2.4333333333333336,3.183333333333333)},x=1.6666666666666667cm,y=0.7666666666666666cm]\draw[->,gray] (-0.8,0)--(2.8,0) node[right]{$x$};\draw[->,gray] (0,-2)--(0,4) node[above]{$y$};\draw[blue,thick,solid,domain=-0.4:2.4,samples=180] plot(\x,{2*(\x-1)^2-1});\fill[orange!85!black] (1,-1) circle(2.2pt);\fill[orange!85!black] (0.2928932188134524,0) circle(2.2pt);\fill[orange!85!black] (1.7071067811865475,0) circle(2.2pt);\draw[violet,dashed] (1,-.95)--(1,3.6);\node[below] at (1,-1){$(1,-1)$};\end{scope}\node[anchor=west,align=left,text width=5.0cm] at (8.4,4.0) {$f=2(x-1)^2-1$\\[9pt]轴 $x=1$\\[9pt]顶点 $(1,-1)$\\[9pt]$x_{1,2}=1\pm\frac{\sqrt2}{2}$\\[9pt]$\Delta=8>0$};
\node[text width=12.6cm,align=center] at (7.1,.65) {两个零点到轴的距离都是根号2除以2；顶点高度为负，开口向上。};
\end{tikzpicture}\end{center}


设 $a(x-h)^2+k=0$，只需移项：
\[
(x-h)^2=-\frac{k}{a}.
\]
若 $-k/a>0$，两个交点横坐标为 $h\pm\sqrt{-k/a}$；
若 $-k/a=0$，只有 $x=h$；若 $-k/a<0$，因为平方不可能是负数，所以没有实数交点。

\medskip

\end{explanationbox}
\clearpage
\begin{explanationbox}
\textbf{\textcolor{CExplanation}{考试里怎么最快选择？}}
\begin{itemize}[leftmargin=*,itemsep=0.3em,topsep=0.3em]
\item 给出顶点 $(h,k)$：先设 $y=a(x-h)^2+k$，通常只剩 $a$ 一个数要找。
\item 给出两个与 $x$ 轴的交点：先设 $y=a(x-x_1)(x-x_2)$，通常也只剩 $a$。
\item 给出三个不同横坐标的点：通常设 $y=ax^2+bx+c$；若其中一个点恰好在 $y$ 轴上，可立即确定 $c$。
\item 只给对称轴 $x=h$：设 $y=a(x-h)^2+k$，\textbf{注意还有 $a,k$ 两个数待求}。
\end{itemize}
\end{explanationbox}
